\documentclass[10pt,firstlineindent,nopartpage,article,accentcolor=tud1b]{tudreport}

\usepackage{booktabs}
\usepackage{amsmath}
\usepackage{float}
\usepackage{hyperref}
\usepackage{enumitem}
\usepackage{siunitx}

\title{Performance Analysis and Optimization of MiniCFD on a Shared-Memory Node}
\SCAuthor{Anonymous Submission}

\begin{document}
\maketitle

\begin{abstract}
This paper studies the OpenMP performance of MiniCFD on a single node of the
Lichtenberg cluster. The focus is the pressure-correction phase, which dominates
the runtime of the application and limits parallel scaling. A reproducible
benchmark campaign was built around an output-free run configuration, automated
Slurm execution, repeated measurements, and HPCToolkit profiling. The baseline
configuration with the default DIC preconditioner showed very weak scaling:
runtime improved only from \SI{136.22}{s} at one thread to \SI{108.52}{s} at
16 threads, and higher thread counts did not help. Profiling identified two
main issues: expensive pressure-solver work inside preconditioned conjugate
gradient (PCG), and unnecessary temporary-vector overhead in the solver loop.
The implemented optimizations were a reduction of shared grid-access overhead in
the advection phase, in-place vector updates in PCG, and reuse of already
computed solver quantities. In the final version, the DIC configuration
improved to \SI{81.66}{s} at 16 threads and \SI{82.30}{s} at 104 threads. A
comparative Jacobi campaign, starting from the common Opt-A code state, showed
that the same solver-side changes reduced the 8-thread runtime further to
\SI{80.00}{s}. The remaining bottleneck is the limited scalability of the
pressure solver, especially the sequential DIC preconditioner and the
memory-bound sparse matrix-vector product.
\end{abstract}

\section{Introduction}
MiniCFD is a compact teaching application for incompressible flow simulation.
Its main value in a performance-engineering course is not peak realism, but the
fact that it contains a complete numerical workflow in a code base that is
small enough to inspect directly. That makes it a good target for a
single-node study: timing can be automated, hotspots can be inspected in
detail, and code changes can be related back to concrete numerical kernels.

The goal of this work is to answer a focused question: what limits the
performance of MiniCFD on one shared-memory node, and which small, justified
optimizations can improve it? The scope is deliberately restricted to OpenMP on
one Lichtenberg node. No MPI scaling study, compiler comparison, or multi-node
analysis is included.

The study follows a standard performance-engineering loop. First, a stable
baseline is established. Second, hotspot analysis is used to identify the
dominant costs. Third, small code changes are introduced one at a time and
re-measured under the same benchmark conditions. This approach keeps causality
clear: every result can be tied back to one specific change or one specific
design decision.

\section{Related Work and Tooling}
The numerical core of MiniCFD is a standard sparse iterative-solver problem.
Preconditioned conjugate gradient and related methods are classical tools for
large sparse linear systems. Saad~\cite[pp.~107,~113]{saad} discusses how
classical iterations such as Jacobi and Gauss-Seidel can be interpreted through
preconditioning, which is exactly the solver setting relevant for MiniCFD.
From that perspective, the MiniCFD pressure solver is representative: it
combines a sparse matrix-vector product, repeated dot products, and
preconditioner application inside a tight convergence loop.

For performance analysis, Adhianto et al.~\cite[p.~685]{hpctoolkit} describe
HPCToolkit as an integrated suite for measurement, analysis, attribution, and
presentation of performance data for sequential and parallel programs. They
also motivate measurement without intrusive instrumentation for optimized code
\cite[p.~686]{hpctoolkit}. OpenMP behavior was interpreted together
with source inspection instead of relying on runtime metrics alone. This is
important for MiniCFD because some limits are not caused by a bad OpenMP
setting, but by the algorithmic structure of the underlying kernels.

\section{Application Background}
MiniCFD evolves a velocity field on a uniform three-dimensional Cartesian grid.
For the experiments in this paper, the command-line parameters have the
following meaning:
\begin{itemize}[leftmargin=*]
\item \texttt{-d}: number of grid cells per spatial axis,
\item \texttt{-e}: simulation end time,
\item \texttt{-s}: time-step size,
\item \texttt{-p}: preconditioner choice for the pressure solver.
\end{itemize}

For a domain size of $d=96$, the simulation grid contains
$96^3=884{,}736$ cells. With \texttt{-e 10} and \texttt{-s 0.4}, the
application performs $10 / 0.4 = 25$ time steps. Each time step consists
mainly of three phases:
\begin{enumerate}[leftmargin=*]
\item semi-Lagrangian advection of the velocity field,
\item pressure correction by solving a sparse linear system,
\item application of the pressure correction to enforce incompressibility.
\end{enumerate}

The pressure correction is the most expensive phase. In the current code, the
pressure matrix is assembled from a 7-point Laplace stencil in every time step.
The resulting sparse system is solved with a preconditioned conjugate gradient
(PCG) method. Two preconditioners were studied:
\begin{itemize}[leftmargin=*]
\item Jacobi, which is simple and cheap per iteration,
\item DIC, a simplified diagonal-based incomplete Cholesky variant.
\end{itemize}

In this code base, this creates a practical tradeoff. A stronger preconditioner
can reduce the number of PCG iterations, but each preconditioner application is
more expensive. A weaker preconditioner is cheaper to apply, but it can leave
more work to the Krylov solver and therefore increase the total cost of sparse
matrix-vector products and reductions over the whole iteration process.

This structure already suggests where bottlenecks may appear. In MiniCFD, the
sparse matrix-vector product performs indirect accesses through the sparse
matrix structure and repeatedly gathers values from the input vector. The DIC
implementation also performs forward and backward sweeps over dependent entries.
Both patterns restrict scalable parallel work, which is consistent with the
observed profiles and scaling curves.

\section{Experimental Setup}
\subsection{Platform and Build}
All main measurements were performed on the Lichtenberg cluster using the
course reservation \texttt{kurs00102}. Jobs were executed through Slurm on a
single node with up to 104 hardware threads. The code was compiled with GCC
\texttt{11.5.0}. The main benchmark build used \texttt{RelWithDebInfo} with
\texttt{-O2 -g -DNDEBUG}. Profiling tools were provided through the prepared
Spack environment.

The work used two storage locations according to the cluster policy. The Git
repository, scripts, and derived summaries stayed in the persistent home
directory. Larger raw outputs, including per-run logs and HPCToolkit data, were
written to \texttt{\$HPC\_SCRATCH}.

\subsection{Benchmark Design}
Benchmark automation was added to the repository in the form of build scripts,
campaign generators, Slurm wrappers, result collectors, and CSV aggregators.
For performance measurements, a new command-line flag
\texttt{--disable-output} was introduced so that file I/O did not distort the
compute-oriented benchmarks.

The pilot campaign compared domain sizes 64, 96, and 128 with the DIC
preconditioner. All runs used \texttt{-s 0.4}, disabled output, and three
repetitions. The results are shown in Table~\ref{tab:pilot}. Domain size 64 was
already measurable, but 96 was chosen as the main benchmark because it offered
a better compromise between runtime, profiling visibility, and cluster queue
constraints. Domain size 128 was unnecessarily expensive for repeated
experiments.

\begin{table}[H]
\centering
\caption{Pilot measurements used to choose a representative problem size.}
\label{tab:pilot}
\begin{tabular}{rrrr}
\toprule
Domain size & Threads & Mean runtime [s] & 16-thread speedup \\
\midrule
64  & 1  & 30.48  & -- \\
64  & 16 & 22.37  & 1.35 \\
96  & 1  & 135.92 & -- \\
96  & 16 & 108.74 & 1.25 \\
128 & 1  & 537.57 & -- \\
128 & 16 & 464.70 & 1.16 \\
\bottomrule
\end{tabular}
\end{table}

The final benchmark setup used:
\begin{itemize}[leftmargin=*]
\item domain size \texttt{-d 96},
\item end time \texttt{-e 10},
\item time-step size \texttt{-s 0.4},
\item disabled output,
\item three repetitions per configuration.
\end{itemize}

The DIC campaigns form the full baseline-to-final optimization chain. The
Jacobi campaigns were added later, after Optimization~A was already present in
the working tree. As a result, the Jacobi progression reported in this paper
starts from the common Opt-A code state rather than from a pre-\texttt{opt-a}
baseline.

An OpenMP placement calibration compared \texttt{OMP\_PROC\_BIND=close} and
\texttt{spread} at 16, 32, 52, and 104 threads. The difference was negligible:
for example, the 16-thread DIC runtime changed only from \SI{108.04}{s}
(\texttt{close}) to \SI{108.72}{s} (\texttt{spread}). Since the performance was
effectively identical, \texttt{spread} was kept for the remaining experiments.

\subsection{Measurement Quality}
Relative standard deviations in the main timing campaign stayed below one
percent in all reported DIC runs and below half a percent in the main Jacobi
runs. For example, the baseline DIC case ranged from
\SI{0.16}{\percent} to \SI{0.71}{\percent}. This indicates that the benchmark
setup was stable enough for comparing small code changes.

\subsection{Automation and Reproducibility}
An important part of the project was not only the measurements themselves, but
the ability to repeat them. The repository was extended with a small automation
layer for cluster builds, campaign generation, Slurm execution, result
collection, aggregation, and plotting. Every run produced machine-readable
metadata, including the Git commit, build name, compiler flags, node name,
thread count, OpenMP placement, simulation parameters, exit code, and wall
time. This made it possible to compare versions without relying on manually
copied command lines or ad hoc notes.

This reproducibility work mattered in practice. Several optimization ideas were
tested in close succession, including one intentionally unsuccessful change.
Without automated campaign generation and aggregated CSV summaries, it would
have been very difficult to separate real improvements from random variation or
configuration mistakes. The measurement pipeline therefore became part of the
engineering result, not just a convenience script around it.

\section{Baseline Characterization}
Table~\ref{tab:baseline-dic} shows the baseline scaling of the default DIC
configuration. The most important result is not the absolute runtime, but the
shape of the scaling curve. Performance improves somewhat up to 8--16 threads,
then essentially stops improving. At 104 threads, the runtime is even slightly
worse than at 16 threads.

\begin{table}[H]
\centering
\caption{Baseline scaling with the default DIC preconditioner.}
\label{tab:baseline-dic}
\begin{tabular}{rrrr}
\toprule
Threads & Mean runtime [s] & Speedup & Efficiency \\
\midrule
1   & 136.22 & 1.00 & 1.000 \\
2   & 124.56 & 1.09 & 0.544 \\
4   & 113.53 & 1.20 & 0.299 \\
8   & 109.68 & 1.24 & 0.155 \\
16  & 108.52 & 1.25 & 0.078 \\
32  & 109.25 & 1.24 & 0.039 \\
52  & 109.31 & 1.24 & 0.024 \\
64  & 109.52 & 1.24 & 0.019 \\
104 & 112.81 & 1.20 & 0.012 \\
\bottomrule
\end{tabular}
\end{table}

This baseline already reveals the core problem: MiniCFD uses OpenMP, but the
pressure-solver path does not expose enough scalable work to benefit from very
high thread counts. A pure placement change cannot solve that. The reason must
be inside the numerical kernels themselves.

\section{Optimizations}
Unless stated otherwise, the optimization stages are cumulative:
Opt-C builds on Opt-A, and Opt-D builds on Opt-C. The only exception is
Opt-B, which was an isolated Jacobi experiment and was not carried forward.

\subsection{Optimization A: Remove Hot-Path Grid-Access Overhead}
The first optimization was motivated by the baseline DIC profiles. The serial
run showed that \texttt{solveAdvection(...)} still consumed 13.4\% of total
runtime, which was too large to ignore early. The 16-thread worker-thread view
then showed a much more specific pathology: 50.2\% of worker-thread time was
spent inside \texttt{semiLagrangianAdvection(...)}, and 24.2\% was spent in
\texttt{gomp\_simple\_barrier\_wait}. Drilling into this region identified
\texttt{traceBackward(...)} and \texttt{VelocityField::trilerp(...)} as the hot
subroutines.

The corresponding source-level issue was small but highly repeated. The
advection code accessed the grid through a \texttt{shared\_ptr} interface in a
hot OpenMP region. A const grid-reference accessor was added and used in the
advection path, so that the numerical method remained unchanged while pointer
and reference-management overhead were removed from the hottest inner loop.

The performance effect was modest in serial but visible in parallel. For the
default DIC configuration, the 16-thread runtime improved from \SI{108.52}{s}
to \SI{103.69}{s}, while the one-thread runtime changed only from
\SI{136.22}{s} to \SI{134.92}{s}. This was exactly the expected pattern for a
parallel hot-path cleanup. Table~\ref{tab:opt-a-effect} summarizes the effect
at representative thread counts.

\begin{table}[H]
\centering
\caption{Effect of Optimization~A on the DIC configuration.}
\label{tab:opt-a-effect}
\begin{tabular}{rrrr}
\toprule
Threads & Baseline [s] & Opt-A [s] & Change [\%] \\
\midrule
1   & 136.22 & 134.92 & -1.0 \\
8   & 109.68 & 104.78 & -4.5 \\
16  & 108.52 & 103.69 & -4.4 \\
104 & 112.81 & 105.32 & -6.6 \\
\bottomrule
\end{tabular}
\end{table}

\subsection{Optimization C: Replace Temporary-Heavy PCG Vector Expressions}
After Optimization~A, the advection phase became less dominant and the next
profile clearly pointed at the solver. In the post-Opt-A DIC profile, the main
costs were concentrated in three places: the PCG driver, the DIC
preconditioner, and the sparse matrix-vector product. The same profile also
showed repeated cost in vector operators such as addition, scalar
multiplication, and assignment.

The serial Opt-A profiles were especially useful because they separated generic
solver overhead from preconditioner-specific cost. In the DIC Opt-A serial run,
\texttt{DICPreconditioner::precondition(...)} still dominated with 34.0\%,
\texttt{spmv(...)} contributed 16.0\%, and
\texttt{solveAdvection(...)} remained at 13.5\%. In the Jacobi Opt-A serial
run, the balance shifted strongly toward generic PCG work:
\texttt{spmv(...)} rose to 29.8\%, the Jacobi preconditioner itself accounted
for only 4.8\%, and several overloaded vector operators inside the same PCG
loop each contributed another 4--8\%. The key conclusion was that the next
useful change should target the shared PCG update loop rather than one specific
preconditioner implementation.

The relevant code pattern inside PCG was
\begin{equation*}
x = x + \alpha d,\qquad
r = r - \alpha z,\qquad
d = h + \beta d.
\end{equation*}
One PCG iteration therefore performs one sparse matrix-vector product, several
dot products, one preconditioner application, and then these three vector
updates. HPCToolkit showed that this loop body was executed often enough that
apparently ``small'' operator-overload costs accumulated into a visible
fraction of total runtime.

This style is compact, but each expression may allocate or copy temporary
vectors. Since the PCG loop executes many times per time step, even moderate
temporary-creation overhead becomes expensive in total.

The change was therefore to replace the temporary-heavy expressions with
in-place vector helpers: one AXPY-style update for the solution, one
AXPY-style update for the residual, and one combined update for the search
direction. This did not alter the solver mathematics. It only reduced
intermediate objects and memory traffic.

This was the largest single optimization of the project. For DIC, the 16-thread
runtime improved from \SI{103.69}{s} to \SI{84.87}{s}. For Jacobi, which was
measured from the same Opt-A code state, the 16-thread runtime improved from
\SI{163.66}{s} to \SI{97.37}{s}. The stronger effect in Jacobi is consistent
with the HPCToolkit profile: after Opt-C, Jacobi spent about 43.1\% of its time
in \texttt{spmv(...)} and 16.5\% in dot products, while
\texttt{JacobiPreconditioner::precondition(...)} itself contributed only 7.9\%.
In other words, Jacobi exposed more of the generic PCG machinery, so reducing
vector overhead helped more. Table~\ref{tab:optac-compare} compares the
preconditioners on matching thread counts before and after this solver rewrite.

\begin{table}[H]
\centering
\caption{Opt-A to Opt-C progression on matching thread counts for both preconditioners.}
\label{tab:optac-compare}
\begin{tabular}{rrrrr}
\toprule
Threads & DIC Opt-A [s] & DIC Opt-C [s] & Jacobi Opt-A [s] & Jacobi Opt-C [s] \\
\midrule
1   & 134.92 & 113.03 & 223.67 & 153.46 \\
8   & 104.78 & 86.06  & 165.37 & 99.33  \\
16  & 103.69 & 84.87  & 163.66 & 97.37  \\
104 & 105.32 & 84.88  & 169.94 & 106.10 \\
\bottomrule
\end{tabular}
\end{table}

\subsection{Optimization D: Reuse Existing Solver Quantities}
After Optimization~C, the remaining solver overhead was smaller, but still
visible. The next low-risk target was repeated work inside the same PCG loop.
Two opportunities were used.

First, the step size originally recomputed an inner product that was already
available as \texttt{oldSqrResidNorm}. Second, Jacobi preconditioning still
returned a fresh vector on every application, even though the solver could
reuse a preallocated output vector. Both issues are small in isolation, but
they occur in the hottest iterative region of the program.

The change was therefore to reuse the existing scalar
\texttt{oldSqrResidNorm} in the step-size calculation and to add an in-place
preconditioner interface for Jacobi. This again preserved the numerical method.
It only removed repeated work and repeated object creation.

The result was a smaller but still measurable improvement. For DIC, the
16-thread runtime improved from \SI{84.87}{s} to \SI{81.66}{s}. For Jacobi, the
16-thread runtime improved from \SI{97.37}{s} to \SI{81.00}{s}, and the
8-thread runtime improved from \SI{99.33}{s} to \SI{79.97}{s}. This confirmed
that the solver still contained profitable low-level inefficiencies even after
the larger Opt-C rewrite.

\subsection{Optimization B: Failed OpenMP Experiment on Jacobi}
One negative result is worth keeping because it clarifies what did not work. A
separate experiment parallelized the Jacobi preconditioner itself with a simple
OpenMP loop. This looked attractive because Jacobi is structurally simpler than
DIC, but the measurement contradicted that intuition. At 8 threads, the Jacobi
runtime worsened from \SI{165.37}{s} to \SI{167.66}{s}, and at 104 threads it
worsened from \SI{169.94}{s} to \SI{178.66}{s}.

This failed optimization supports the main argument of the paper: the project
was not improved by adding directives blindly. Improvements came from following
the measured hotspot sequence. In this case, the full PCG path mattered more
than a naive OpenMP change to one small component.

\subsection{Performance Summary}
The cumulative DIC progression is shown in Table~\ref{tab:dic-progression}.
Optimization A is mainly a parallel cleanup. Optimization C delivers the main
solver improvement. Optimization D adds a smaller final refinement.

\begin{table}[H]
\centering
\caption{Runtime progression for the default DIC configuration.}
\label{tab:dic-progression}
\begin{tabular}{lrrr}
\toprule
Version & 1 thread [s] & 16 threads [s] & 104 threads [s] \\
\midrule
Baseline & 136.22 & 108.52 & 112.81 \\
Opt-A    & 134.92 & 103.69 & 105.32 \\
Opt-C    & 113.03 & 84.87  & 84.88 \\
Opt-D    & 110.70 & 81.66  & 82.30 \\
\bottomrule
\end{tabular}
\end{table}

Compared to baseline, the final DIC version reduced runtime by about
\SI{18.7}{\percent} at one thread, \SI{24.8}{\percent} at 16 threads, and
\SI{27.0}{\percent} at 104 threads. The absolute improvement is therefore real
and substantial. However, the scaling curve is still unsatisfactory. Even in
the final version, the speedup at 16 threads is only 1.35, and performance
stays almost flat between 16 and 104 threads.

\begin{table}[H]
\centering
\caption{Final Opt-D comparison between the two preconditioners at matching thread counts.}
\label{tab:final-compare}
\begin{tabular}{rrr}
\toprule
Threads & DIC [s] & Jacobi [s] \\
\midrule
1   & 110.70 & 135.06 \\
2   & 97.63  & 103.25 \\
4   & 88.42  & 86.64 \\
8   & 83.93  & 79.97 \\
16  & 81.66  & 81.00 \\
32  & 81.78  & 87.69 \\
52  & 80.93  & 87.24 \\
104 & 82.30  & 92.03 \\
\bottomrule
\end{tabular}
\end{table}

The final comparison in Table~\ref{tab:final-compare} is instructive because
both columns come from the same final Opt-D code state and therefore differ
only in the selected preconditioner. DIC is clearly better in serial execution
because its stronger preconditioner reduces solver work more effectively.
Jacobi, however, becomes competitive once the solver-side overheads are
reduced. At 8 and 16 threads, the two methods reach almost identical runtimes,
even though they arrive there for different reasons: DIC spends more time in
preconditioning, while Jacobi shifts more work toward \texttt{spmv()} and
reductions. This supports the main profiling conclusion that MiniCFD is limited
less by one universal bottleneck and more by the interaction between solver
structure and memory behavior.

The same table also makes the scaling limit visible in a way that a single
curve does not. Both preconditioners improve clearly from 1 to 8 threads.
Between 8 and 16 threads, the gains become small. Beyond that point the curves
either flatten completely (DIC) or turn upward again (Jacobi). This is the
measured consequence of the hotspot analysis: once the code is dominated by
memory-bound \texttt{spmv()}, reductions, and in the DIC case additionally by
sequential triangular sweeps, adding more OpenMP threads no longer creates
proportional useful work.

\subsection{Why Scaling Still Flattens}
The final code is clearly faster than the baseline, but it still does not scale
like a highly optimized production solver. There are several reasons:
\begin{itemize}[leftmargin=*]
\item The DIC preconditioner contains sequential forward and backward sweeps.
\item The sparse matrix-vector product uses indirect memory accesses and is
therefore bandwidth-sensitive.
\item The PCG loop is mathematically sequential across iterations. Only the
vector kernels inside each iteration can be parallelized.
\end{itemize}

These points also explain why thread counts above 8 or 16 bring little benefit.
The application does not have enough scalable work per iteration to keep 32, 52,
or 104 threads productive. In the profiles, this appears as substantial OpenMP
waiting time.

\section{Discussion}
The most useful lesson from this study is that performance improvements did not
come from one large algorithmic rewrite. They came from identifying where the
program actually spends time and then removing overheads in the hottest path.

The strongest implemented optimization was not a new numerical method. It was a
change in data movement: replacing temporary-heavy vector expressions with
in-place updates. Even when the floating-point formulas stay identical, the
cost of allocation, copying, and extra passes over memory can dominate a large
fraction of runtime.

The second lesson is methodological. The failed Jacobi OpenMP experiment was
not wasted work. It ruled out an intuitive but incorrect optimization strategy
and made the later improvements more defensible. In a performance paper,
negative results are useful when they narrow the search space and support the
final argument.

Finally, the profiles show that MiniCFD still contains further optimization
potential. The most important remaining limitation is the pressure solver
itself. Beyond the implemented low-level cleanups, the code would benefit from
a more parallel-friendly solver and preconditioner combination than the current
PCG plus DIC path.

\section{Conclusion}
This paper presented a single-node OpenMP performance study of MiniCFD on
Lichtenberg. The baseline application showed weak scaling with the default DIC
preconditioner: 16 threads improved runtime only from \SI{136.22}{s} to
\SI{108.52}{s}. HPCToolkit identified the pressure-correction phase as the main
bottleneck, with different dominant kernels for DIC and Jacobi.

Three useful optimization directions were validated. First, repeated hot-path
grid access in advection was reduced. Second, the PCG loop was rewritten to use
in-place vector updates instead of temporary-heavy expressions. Third, already
computed solver quantities were reused and Jacobi preconditioning was converted
to an in-place form. Together, these changes reduced the DIC runtime to
\SI{81.66}{s} at 16 threads and the Jacobi runtime to \SI{79.97}{s} at 8
threads.

The remaining limitation is not measurement noise or poor job placement. It is
algorithmic and structural: DIC contains sequential sweeps, \texttt{spmv()} is
memory-bound, and the PCG loop exposes only limited thread-parallel work per
iteration.

\section{Future Work}
The most relevant future work is not more micro-tuning of OpenMP settings, but
reconsidering the pressure-solver design itself.

The pressure solver would benefit from a more parallel-friendly algorithmic
combination, for example a different preconditioner or a solver variant that
exposes more useful work inside each iteration.

\section*{Disclosure}
OpenAI Codex was used for wording refinement and \LaTeX{} editing support.
Benchmark data, profiling runs, code execution, interpretation of results, and
the final judgment are the responsibility of the author.

\begin{thebibliography}{9}

\bibitem{saad}
Y.~Saad.
\newblock \emph{Iterative Methods for Sparse Linear Systems}.
\newblock Society for Industrial and Applied Mathematics, 2nd edition, 2003.

\bibitem{hpctoolkit}
S.~Adhianto, S.~Banerjee, M.~Fagan, M.~Kruskal, Q.~Liu, L.~Yuan, and
W.~N. Scherer.
\newblock HPCToolkit: Tools for performance analysis of optimized parallel
programs.
\newblock \emph{Concurrency and Computation: Practice and Experience},
22(6):685--701, 2010.

\end{thebibliography}

\end{document}
